\chapter{Wie entsteht ein Held?}

Jeder Held beginnt seine Reise als Abenteurer.

Bevor ein Held sein erstes Abenteuer erlebt, erhält er eine Geschichte, Werte, Fähigkeiten und eine passende Ausrüstung. Diese Entscheidungen bestimmen, wer der Held ist, worin er besonders gut ist und auf welche Weise er Herausforderungen meistert.

Im Laufe seiner Abenteuer sammelt ein Held Erfahrung. Dadurch verbessert er seine Fähigkeiten, erlernt neue Talente und Zauber und entwickelt sich vom Abenteurer über den Helden bis hin zu einer Legende.

Dieses Kapitel beschreibt sowohl die Erschaffung eines neuen Helden als auch seine Weiterentwicklung während einer Kampagne.

\section{Einen Helden erschaffen}

Die Heldenerschaffung erfolgt in mehreren Schritten. Die folgende Checkliste führt durch den gesamten Ablauf.

\begin{kurzregel}[title=Checkliste: Einen Helden erschaffen]
\begin{itemize}[label=$\square$]
  \item Hintergrund und Geschichte ausdenken
  \item Volk wählen
  \item Rolle wählen
  \item Attribute verteilen
  \item Fertigkeiten verteilen
  \item Abgeleitete Werte berechnen
  \item Ein Starttalent wählen
  \item Startzauber wählen
  \item Passende Startausrüstung auswählen
  \item Persönlichkeit und Aussehen ausgestalten
\end{itemize}
\end{kurzregel}

Sind alle Schritte abgeschlossen, ist der Held bereit für sein erstes Abenteuer.

\subsection{Hintergrund und Geschichte ausdenken}

Überlege zunächst, wer dein Held ist und welches Leben er geführt hat, bevor sein erstes Abenteuer beginnt.

Woher kommt er? Was hat er bisher getan? Warum zieht er ins Abenteuer? Welche Ziele verfolgt er, und welche Menschen, Orte oder Ereignisse haben ihn geprägt?

Der Hintergrund muss nicht bis ins letzte Detail feststehen. Einige Antworten können sich erst während des Spiels entwickeln. Er sollte jedoch eine Vorstellung davon vermitteln, wer der Held ist und warum seine späteren Entscheidungen zu ihm passen.

\subsection{Volk wählen}

Wähle, welchem Volk dein Held angehört. Welche Völker in einem Abenteuer vorkommen und welche Bedeutung sie in der Spielwelt besitzen, legen die Spielgruppe und die Spielleitung gemeinsam fest.

Die Wahl des Volkes ist Teil der Geschichte des Helden. Sie kann seine Herkunft, seine Erfahrungen und die Art beeinflussen, wie andere Figuren ihm begegnen.

\subsection{Rolle wählen}

Wähle eine Rolle, die beschreibt, auf welche Weise dein Held seiner Gruppe zu Beginn besonders helfen möchte: Draufgänger, Verteidiger, Anführer oder Taktiker.

Die Rolle ist weder ein Beruf noch eine feste Klasse. Sie dient als Ausgangspunkt für die Entwicklung des Helden. Später kann er Rollentalente aus verschiedenen Rollen erlernen und miteinander kombinieren.

Die Rollen und ihre Schwerpunkte werden in Kapitel 7 -- Was sind Talente? beschrieben.

\subsection{Attribute verteilen}

Zu Beginn der Heldenerschaffung erhält der Spieler acht grüne Attributswürfel, die frei auf die acht Attribute verteilt werden können.

\begin{kurzregel}[title=Attribute verteilen]
Verteile acht grüne Attributswürfel auf die Attribute. Kein Attribut darf zu Beginn mehr als zwei Würfel besitzen. Attribute dürfen auch keinen Würfel besitzen.
\end{kurzregel}

Durch die Verteilung legt der Spieler die natürlichen Stärken und Schwächen seines Helden fest.

\begin{beispiel}[title=Kieras Attribute]
Kiera soll mutig, aufmerksam, magisch begabt und beweglich sein. Sie investiert jeweils zwei Würfel in MUT und INT sowie jeweils einen Würfel in CHA, GEW, KON und KRA.

Damit besitzt sie MUT 2, KLU 0, INT 2, CHA 1, FIN 0, GEW 1, KON 1 und KRA 1. Insgesamt hat sie alle acht Attributswürfel verteilt.
\end{beispiel}

\subsection{Fertigkeiten verteilen}

Zu Beginn der Heldenerschaffung erhält der Spieler neun blaue Fertigkeitswürfel, die frei auf die Fertigkeiten verteilt werden.

\begin{kurzregel}[title=Fertigkeiten verteilen]
Verteile neun blaue Fertigkeitswürfel auf die Fertigkeiten. Keine Fertigkeit darf zu Beginn mehr als zwei Würfel besitzen. Fertigkeiten dürfen auch keinen Würfel besitzen.
\end{kurzregel}

Die Fertigkeiten bestimmen, in welchen Bereichen ein Held bereits Erfahrung gesammelt hat. Eine Übersicht aller Fertigkeiten befindet sich in Kapitel 2 -- Was macht einen Helden aus? Die Regeln für Fertigkeitsproben stehen in Kapitel 3 -- Wie funktionieren Proben?

\subsection{Abgeleitete Werte berechnen}

Nachdem die Attribute verteilt wurden, werden die davon abgeleiteten Werte bestimmt.

\begin{kurzregel}[title=Abgeleitete Werte]
\textbf{Lebenspunkte} $= 14 + 2 \times \mathrm{KON}$

\textbf{Manavorrat} $= 2 \times \mathrm{CHA}$
\end{kurzregel}

Talente und andere Effekte können diese Werte anschließend verändern. Die genauen Regeln für Lebenspunkte und Mana stehen in Kapitel 2 -- Was macht einen Helden aus?

\subsection{Ein Starttalent wählen}

Jeder Held beginnt das Spiel mit einem Starttalent. Als Starttalent darf ein beliebiges Talent gewählt werden, dessen Voraussetzungen der Held beim Spielstart erfüllt.

Das Starttalent gilt als erlernt und trainiert. Die Regeln für Talente sowie die Bedeutung der Rollen stehen in Kapitel 7 -- Was sind Talente?

\subsection{Startzauber wählen}

Hat ein Held das Talent \emph{Magiebegabt} gewählt, kann er bereits zu Beginn Zauber beherrschen.

\begin{kurzregel}[title=Startzauber]
Die Summe der Manakosten aller Startzauber darf den Manavorrat des Helden nicht überschreiten. Kein einzelner Startzauber darf mehr als 4 Mana kosten. Jeder Zauber kann nur einmal erlernt werden.
\end{kurzregel}

Das Talent \emph{Magiebegabt} erhöht den Manavorrat des Helden um 6 Mana. Die Regeln für Zauber stehen in Kapitel 6 -- Wie wirkt Magie?

\subsection{Passende Startausrüstung auswählen}

Jeder Held beginnt das Spiel mit einer Ausrüstung, die zu seiner bisherigen Geschichte passt.

Der Spieler stellt die Startausrüstung gemeinsam mit der Spielleitung zusammen. Sie sollte zu einem neuen Helden auf dem Rang Abenteurer sowie zu seinem sozialen und finanziellen Hintergrund passen. Die Spielleitung entscheidet, welche Gegenstände für den Helden glaubwürdig und angemessen sind.

Ein Jäger könnte beispielsweise einen Bogen und leichte Kleidung besitzen. Eine ehemalige Wache könnte mit Schwert, Schild und Rüstung beginnen. Ein wandernder Heiler trägt vielleicht einen Stab, einen Beutel mit Kräutern und einfache Reisekleidung.

\subsection{Persönlichkeit und Aussehen ausgestalten}

Zum Abschluss erhält der Held die Merkmale, durch die er am Spieltisch lebendig wird. Dazu können unter anderem gehören:

\begin{itemize}
  \item Name und Alter,
  \item Haar- und Augenfarbe,
  \item Körperbau und auffällige Merkmale,
  \item Kleidung und persönlicher Stil,
  \item Eigenheiten, Vorlieben und Abneigungen,
  \item Beziehungen zu anderen Figuren sowie
  \item Hoffnungen, Ängste und persönliche Ziele.
\end{itemize}

Nicht jedes Detail muss bereits vor dem ersten Abenteuer feststehen. Die Figur darf sich während des gemeinsamen Spiels weiterentwickeln.

\section{Beispiel: Kiera erschaffen}

Das folgende Beispiel führt die Heldenerschaffung einmal vollständig durch.

\begin{beispiel}[title=Kiera erschaffen]
\textbf{Hintergrund:} Kiera wuchs als Tochter einer einfachen Handwerkerfamilie auf und arbeitete einige Jahre als Botin. Auf ihren Wegen lernte sie, Gefahren früh zu erkennen, sich unauffällig zu bewegen und notfalls mit einem Speer zu verteidigen. Seit sie erstmals unkontrolliert Magie gewirkt hat, sucht sie nach Antworten über die Herkunft ihrer Kräfte.

\textbf{Volk:} Kiera ist ein Mensch.

\textbf{Rolle:} Sie beginnt als Draufgängerin. Kiera möchte Gegner entschlossen angreifen, ist aber nicht dauerhaft an diese Rolle gebunden.

\textbf{Attribute:} Kiera verteilt ihre acht Attributswürfel auf MUT 2, INT 2, CHA 1, GEW 1, KON 1 und KRA 1. KLU und FIN bleiben auf 0.

\textbf{Fertigkeiten:} Sie verteilt ihre neun Fertigkeitswürfel auf Schlagen 2, Zaubern 2, Schleichen 2, Turnen 1, Suchen 1 und Überreden 1. Alle übrigen Fertigkeiten bleiben auf 0.

\textbf{Abgeleitete Werte:} Mit KON 1 besitzt Kiera $14 + 2 \times 1 = 16$ Lebenspunkte. Mit CHA 1 besitzt sie zunächst 2 Mana.

\textbf{Starttalent:} Kiera wählt \emph{Magiebegabt}. Das Talent erhöht ihren Manavorrat um 6, sodass sie mit 8 Mana beginnt.

\textbf{Startzauber:} Kiera erlernt \emph{Weißer Speer} für 1 Mana, \emph{Zielscheibe} für 2 Mana, \emph{Geisterschritt} für 3 Mana und \emph{Herzschlag spüren} für 2 Mana. Die Summe ihrer Manakosten entspricht genau ihrem Manavorrat von 8.

\textbf{Startausrüstung:} Gemeinsam mit der Spielleitung wählt Kiera einen Speer und eine Knochenrüstung. Diese Ausrüstung passt zu ihrer bisherigen Geschichte und bereitet sie auf ihr erstes Abenteuer vor.

\textbf{Persönlichkeit und Aussehen:} Kiera ist neugierig, direkt und ungeduldig. Sie trägt ihr dunkles Haar meist zusammengebunden, bevorzugt praktische Kleidung und bewahrt eine alte Messingbrosche ihrer Mutter auf. Sie möchte lernen, ihre Magie zu beherrschen, und fürchtet zugleich, durch sie Menschen zu gefährden, die ihr nahestehen.
\end{beispiel}

\section{Einen Helden weiterentwickeln}

Während seiner Abenteuer sammelt ein Held Erfahrung. Mit dieser Erfahrung kann er seine Fähigkeiten verbessern, neue Talente und Zauber erlernen und sich zu einem mächtigeren Helden entwickeln.

\subsection{Erfahrungspunkte erhalten}

Die Spielleitung vergibt Erfahrungspunkte für bestandene Abenteuer, besondere Leistungen oder andere bedeutende Entwicklungen.

Wie genau Erfahrungspunkte vergeben werden, entscheidet jede Spielleitung. Für eine durchschnittliche Kampagne wird empfohlen:

\begin{itemize}
  \item \textbf{Gruppen-EP:} Alle Helden erhalten gleichzeitig die gleiche Menge an EP.
  \item \textbf{EP-Menge:} Etwa 1 bis 2 EP pro Spielabend.
\end{itemize}

\subsection{Heldenränge}

Mit zunehmender Erfahrung steigt ein Held im Rang auf. Es gibt drei Heldenränge: Abenteurer, Held und Legende.

Für den Heldenrang zählt die gesamte Erfahrung, die ein Held im Laufe seiner Abenteuer erhalten hat. Bereits ausgegebene Erfahrungspunkte werden weiterhin mitgezählt.

Jeder Heldenrang erhöht die maximal möglichen Attributs- und Fertigkeitswerte.

\begin{hptable}{L{2.5cm} C{3.0cm} C{3.1cm} C{3.1cm}}
\textbf{Heldenrang} & \textbf{Ab erhaltenen EP} & \textbf{Max. Attributswert} & \textbf{Max. Fertigkeitswert} \\
Abenteurer & 0 & 3 & 3 \\
Held & 75 & 4 & 4 \\
Legende & 150 & 5 & 5 \\
\end{hptable}

Ein Held muss den entsprechenden Heldenrang erreicht haben, bevor er ein Attribut oder eine Fertigkeit über den bisherigen Höchstwert hinaus steigern darf.

\subsection{Erfahrungspunkte ausgeben}

Ein Held kann seine gesammelten Erfahrungspunkte verwenden, um sich weiterzuentwickeln.

\begin{hptable}{Y L{4.5cm}}
\textbf{Verbesserung} & \textbf{Kosten} \\
Fertigkeit steigern & Neuer Fertigkeitswert $\times$ 1 EP \\
Attribut steigern & Neuer Attributswert $\times$ 3 EP \\
Talent erlernen & siehe unten \\
Zauber erlernen & Manakosten $\times$ 2 EP \\
\end{hptable}

Die Kosten zum Erlernen neuer Talente richten sich danach, wie viele Talente der Held bereits erlernt hat.

\begin{hptable}{Y L{4.5cm}}
\textbf{Erlerntes Talent} & \textbf{Kosten} \\
Starttalent & 0 EP \\
2. Talent & 5 EP \\
3. Talent & 10 EP \\
4. Talent & 15 EP \\
Jedes weitere Talent & 15 EP \\
\end{hptable}

Neue Talente gelten sofort als trainiert. Sind dadurch bereits vier Talente derselben Talentart trainiert, entscheidet der Spieler, welches seiner bisher trainierten Talente dieser Talentart stattdessen nicht mehr trainiert ist.